% Generated draft: requires mathematical and source review.
\documentclass[11pt]{article}
\usepackage[margin=25mm]{geometry}
\usepackage{fontspec}
\setmainfont{Latin Modern Roman}
\usepackage{amsmath,amssymb,mathtools,unicode-math}
\setmathfont{Latin Modern Math}
\usepackage{enumitem,needspace,xurl,hyperref}
\setlength{\parindent}{0pt}
\setlength{\parskip}{5pt}
\setlength{\emergencystretch}{3em}
\newcommand{\N}{\mathbb{N}}
\newcommand{\Z}{\mathbb{Z}}
\newcommand{\Q}{\mathbb{Q}}
\newcommand{\R}{\mathbb{R}}
\newcommand{\C}{\mathbb{C}}
\newcommand{\F}{\mathbb{F}}
\newcommand{\A}{\mathbb{A}}
\newcommand{\PP}{\mathbb{P}}
\newcommand{\E}{\mathbb{E}}
\newcommand{\eps}{\varepsilon}
\newcommand{\dd}{\,\mathrm{d}}
\begin{document}
\section*{UnsolvedMath ID: 20002688}
\textbf{Source problem code:} catalogue entry
\subsection{Definitions and mathematical statement}
The source asks whether a diffuse diagonal \(D\), uniform on \([0,1]\), belongs to the von Neumann algebra generated by a band-limited circular operator, and whether that algebra is an interpolated free group factor with a specified parameter. The extraction corrupts the band parameter and freeness symbol. The supplied verification reports a star-freeness assertion and parameter
\[
 r_\alpha=1+2\alpha(1-\alpha/2)=1+2\alpha-\alpha^2.
\]
Use \(\alpha\) for bandwidth, separating it from the circular element denoted by \(c\) in the extraction.

Two readings must be distinguished. If \(D\) is literally scalar-free from \(\{B_\alpha,B_\alpha^*\}\), then \(D\notin N_\alpha=W^*(B_\alpha)\): membership and freeness would force its centered variance to vanish, whereas \(\tau((D-1/2)^2)=1/12\). This conditional observation does not settle a band model in which that freeness premise is false.

The supplied aid interprets the intended model as the Dykema--Tucci characteristic upper-band operator. An explicit convention is a finite tracial algebra containing \(A=L^\infty[0,1]\), \(D(x)=x\), and an \(A\)-valued circular operator \(B_\alpha\), \(0<\alpha\le1\), with conditional covariances
\[
 E_A(B_\alpha f B_\alpha^*)(x)
       =\int_x^{\min(1,x+\alpha)}f(y)\,\dd y,
 \qquad
 E_A(B_\alpha^* f B_\alpha)(x)
       =\int_{\max(0,x-\alpha)}^x f(y)\,\dd y,
\]
for \(f\in A\); the other second circular covariances vanish. The higher conditional moments are those of this operator-valued circular law. This fixes the characteristic upper-band orientation and normalization, which are absent from the damaged source text.

\textbf{Intended questions.} For a proper band \(0<\alpha<1\), is \(D\in N_\alpha\)? Is \(N_\alpha\) isomorphic to \(L(\mathbb F_{r_\alpha})\)? These are not logically exclusive alternatives; the source's word or does not make one the negation of the other. The aid does not resolve either proper-band question.

The aid's exact mixed-moment obstruction in this model is
\[
 \tau\bigl((D-1/2)(B_\alpha^*B_\alpha-
                       \tau(B_\alpha^*B_\alpha))\bigr)
       =\frac{\alpha^2(3-2\alpha)}{12},
\]
while replacing \(B_\alpha^*B_\alpha\) by \(B_\alpha B_\alpha^*\) gives the negative of that value. Hence the literal scalar-freeness premise is incompatible with this intended upper-band construction. For scale, the centered \(B_\alpha^*B_\alpha\) variance in this normalization is
\[
 \tau(B_\alpha^*B_\alpha)=\alpha-\alpha^2/2,\qquad
 \|B_\alpha^*B_\alpha-\tau(B_\alpha^*B_\alpha)\|_2^2
       =\alpha^2(1-\alpha/3-\alpha^2/4).
\]
Writing \(E_{N_\alpha}\) for trace-preserving conditional expectation, the aid reports
\[
 \|E_{N_\alpha}(D)-1/2\|_2
 \ge\frac{\alpha(3-2\alpha)}
 {12\sqrt{1-\alpha/3-\alpha^2/4}}>0.
\]
This proves nontrivial diagonal information is visible in \(N_\alpha\), not full membership. The exact membership test is
\[
 D\in N_\alpha\quad\Longleftrightarrow\quad
 \|E_{N_\alpha}(D)-1/2\|_2^2=1/12,
\]
since conditional expectation is orthogonal projection. In particular a positive lower bound strictly below full variance is not a recovery theorem.

At the full-width endpoint \(\alpha=1\), the aid reports the known quasinilpotent DT case \(W^*(B_1)\cong L(\mathbb F_2)\), consistent with \(r_1=2\). This endpoint does not answer the proper-band questions. The next heading, ``Infinite Divisibility'', attributed to Nica, is adjacent context and supplies no additional question here.
\subsection{Short English statement}
The source's stated freeness and its intended upper-band model are incompatible. Literal scalar freeness excludes diagonal membership, whereas the aid's exact band-model moments show nonzero diagonal correlations. A conditional-expectation norm gives an exact recovery criterion, but the proper-band inclusion and factor-isomorphism questions remain unresolved in the supplied attempt.
\subsection{Sources}
\begin{itemize}
\item Local UnsolvedMath registry, permanent ID 20002688, catalogue entry; canonical \texttt{aim-probability-notes.json}, notes[129].
\item AIM \emph{Free Analysis}, Question 28, page 6, \url{https://aimath.org/WWN/freeanalysis/freeanalysis.pdf}. The supplied verification reports checking the PDF font encoding to recover star freeness and the parameter formula; this draft does not independently repeat that check.
\item The supplied aid identifies the intended upper-band model with Dykema and Tucci and names the full-width quasinilpotent DT case. No paper titles, dates, theorem numbers or external bibliographic URLs are supplied.
\item Adjacent heading ``Infinite Divisibility'', Nica, retained as context only.
\item Machine-generated research attempt 1, partial result, attempt status unknown, reported rigor proved and low novelty confidence; registry status partially solved. Its quantitative recovery reduction is not peer reviewed and requires independent verification.
\end{itemize}
\end{document}
